\chapter{Discussion}
\label{chap:discussion}

\section{Interpretation}

The contribution is an integrative authorization/refusal boundary. Established mechanisms---typed schemas, planning constraints, access policy, task DAGs, leases, fencing, idempotency, and protocol adapters---are combined so that linguistic output cannot silently acquire executable authority. The new contribution claimed here is their operational integration and traceable refinement into one prototype, not invention of each mechanism or formal verification of arbitrary agents.

The LLM is useful where wording varies and decomposition is not known in advance. It is not an authority. Stable capability identifiers, local trust assignment, typed constraints, validation contracts, and completion contracts determine what may be dispatched. A structured rule-only oracle remains a useful upper bound because it shows what can be achieved when interpretation is already supplied; it is not a fair replacement for an end-user system receiving natural language.

Version 0.3 exposed the initial integration's central weakness: safety was obtained through universal refusal. Version 0.4 narrowed the linguistic role and produced substantially better but model-sensitive development results. Version 0.5 then showed that the historical identifier-copying bridge did not transfer to fresh matched pairs: it retained zero unsafe acceptance only by accepting five of 360 feasible repeated observations. That negative experiment motivated the production redesign and is no longer treated as the primary comparison.

The frozen v0.7 study evaluates raw language through production semantic admission and the shared compiler. At primary seed 11, the hybrid attains 71.88\% balanced accuracy and 44.79\% feasible recall, with one unsafe acceptance in 192 observations. The direct LLM attains 53.13\% balanced accuracy and 98.96\% feasible recall but makes 89 unsafe acceptances. The pair-cluster interval for the hybrid-minus-direct unsafe-acceptance rate is entirely negative, so the safety-superiority gate passes. The recall difference is between $-63.54$ and $-45.83$ percentage points, so non-inferiority fails. Safety and utility remain in tension rather than being jointly achieved.

Model separation is essential. Qwen3-1.7B hybrid recall is 56.25\% with no unsafe acceptance; Gemma-4-E2B refuses every feasible case and therefore has zero recall; Qwen3-8B reaches 78.13\% recall but produces the sole unsafe hybrid acceptance. High seed agreement---98.44\% for the hybrid and 96.88\% for direct---shows that the primary pattern is not explained by widespread seed instability. It does not establish backend seed compliance or generalize beyond these quantizations and this laptop runtime.

The exact-name/alias baseline solves all 64 designed cases. Its perfect balanced accuracy means the lexical-superiority gate fails decisively: the hybrid-minus-alias interval is $[-32.81,-23.44]$ percentage points. This is both a negative result for the proposal on the corpus and a construct warning. The designed language contains enough exact aliases for a deterministic matcher to recover the author labels; therefore the corpus does not establish that an LLM is necessary. Future evaluation must include independently sourced utterances whose semantics cannot be recovered through catalog phrase overlap while keeping the same safety boundary.

The production redesign nevertheless isolates where failure occurs. The oracle/compiler invariants all pass, provider selection is registry-order independent, alternative recovery succeeds, and search beyond 4,096 assignments refuses reproducibly. Eight observed primary-seed hybrid acceptances replay through the fixture A2A execution boundary with dependency artifacts and validation contracts, all completing. These results support the implemented deterministic boundary and execution path, not comparative superiority of semantic admission.

\section{Guarantees and Non-Guarantees}

Within the tested implementation, authorization precedes dispatch, unresolved constraints fail closed, declared failed dependencies block descendants, independent branches may continue, session traces are isolated, and stale fenced writes are rejected. These are implementation properties supported by tests and measured traces. They are not mathematical proofs over every deployment, scheduler, transport, or database failure.

Contract validity is weaker than semantic truth. JSON Schema can establish nested type and required-field conformance; a validator can establish a domain-specific property only if that property is encoded. Neither implies that an open-ended summary is complete or that a physical action is safe. The trusted registry and honest-agent assumptions remain explicit.

The admission boundary controls allowed origins, HTTPS policy, locally assigned trust, required credential schemes, and request-header propagation. It does not provide cryptographic identity, attestation, or defense against a malicious admitted agent. Plain HTTP is enabled only for explicit development fixtures.

The legacy PostgreSQL lease-and-fencing mode is not consensus because one database remains its serialization point. The new etcd mode is consensus-backed: Raft orders registry, ownership, plan, attempt, and terminal mutations while a voter majority is available. This does not imply Byzantine tolerance or exactly-once external effects. Fencing prevents an obsolete coordinator from committing authoritative state only when every mutation checks the fence, and external safety additionally requires the receiving agent to enforce the highest fence and durably deduplicate operation keys.

\section{Critical Assessment of the Distributed Implementation}

The strongest design decision is the separation between logical authority and physical availability. The theory still has one accountable authorization role, while etcd removes the single volatile coordinator and single PostgreSQL coordination database from the authoritative path. Transactional owner comparisons, monotonic revision fences, attempt-before-result checks, terminal uniqueness, lease-backed agent registration, and fail-closed quorum errors reinforce one another instead of treating leader election alone as sufficient. Learner-first admission and serialized membership changes reduce reconfiguration risk. Keeping JSONL and PostgreSQL adapters also permits contract comparison and incremental migration rather than an irreversible cutover.

The unified container is operationally simpler, but it creates correlated resource and failure domains. A CPU-intensive planner, slow disk, container restart, host loss, or network namespace fault can affect both the coordinator and its local etcd member. The dispatch concurrency limit mitigates but does not eliminate heartbeat starvation. A production topology must place voters on independent hosts or failure domains, reserve CPU and durable storage for etcd, and monitor disk latency. Three containers on one Docker host demonstrate process-level behavior, not host-level fault tolerance.

Automatic scaling is also easy to overstate. The implementation automatically discovers an existing cluster and reconciles the membership of coordinator nodes that already exist. It does not create compute instances, repair a permanently lost majority, or safely choose a new cluster after loss of quorum. DNS/direct seeds remain necessary across routed networks; multicast is a link-local fallback and may be filtered by container, cloud, or enterprise networks. Target one has no fault tolerance. Targets five and seven tolerate more failures but add write latency and have not yet been exercised by the accepted harness.

The security boundary is appropriate for a trusted private prototype but weak for hostile deployment. HMAC authenticates possession of a cluster-wide secret, not an individual node identity; compromise of one coordinator therefore permits forged cluster messages until the shared secret is rotated. Nonce and timestamp checks depend on bounded clock skew and local replay caches. Discovery authentication does not encrypt or authenticate etcd peer and client transport. Deferring mTLS, certificate rotation, per-node authorization, and attestation leaves etcd ports and shared secrets as high-value assets. The separate agent-registration secret is useful compartmentalization, but admitted malicious agents and compromised voters remain out of scope.

The custom HTTP-gateway client keeps dependencies small and supports endpoint rotation, transactions, leases, and watch resumption, but it also places subtle retry, compaction, cancellation, and error-classification behavior in project code. Quorum-unavailable errors must never be converted into permissive retries or stale reads. The asynchronous PostgreSQL projector correctly avoids a readiness dependency, yet event deletion must remain disabled until projection or export is acknowledged and a retention policy is explicit.

Consensus-v4 shows a materially stronger but still bounded distributed result. Formation at three, five, and seven voters; leader loss; minority partition; majority loss and restoration; concurrent ownership; audit failure; failed-voter replacement; all three crash windows; and the supplementary combined-fault scenario passed in every completed repetition. No executed primary safety invariant failed. Nevertheless, only 140 of 150 primary checks executed because one repetition of each live-reconfiguration path timed out before its five checks ran. These are infrastructure-indeterminate outcomes, not demonstrated safety violations. The campaign also records a dirty source tree and is therefore ineligible as accepted clean evidence. RQ5 is partially supported for the repeatedly completed conditions, while complete reconfiguration reliability and clean-source reproducibility remain open.

\section{Claim and Terminology Discipline}

The writing must distinguish three concepts that earlier drafts conflated. \emph{Centralized authorization} describes one logical decision role; it does not necessarily mean one process. \emph{Consensus-backed replication} describes how that role's state is ordered and recovered; it does not make authorization a peer negotiation. \emph{Automatic membership reconciliation} describes adding or removing already running nodes; it is not automatic compute provisioning. Accordingly, ``fault tolerant'' should always be qualified as crash-fault tolerant while a majority survives, and ``no single point of failure'' should be limited to the authoritative coordination path under independent voter failure domains. The implementation claims duplicate suppression for measured cases, not exactly-once effects.

\section{Threats to Validity}

Construct validity is protected by separating semantic fields, feasibility, execution, artifacts, trace, and recovery measures, but the author-designed and author-labelled corpus creates construct circularity and researcher-bias risk. Freezing labels before confirmatory collection and hiding them during inference reduces leakage, not bias or lack of independent adjudication. Internal validity remains threatened by model nondeterminism, unverifiable backend seed compliance, prompt sensitivity, baseline choices, fixture ownership, one separately classified model-runtime failure, and distributed checks authored with the implementation. External validity is limited to 64 held-out designed cases, three small local models, controlled fixture agents, selected PostgreSQL crash faults, and single-host Docker campaigns. Conclusion validity uses 32 matched pairs, model/category reporting, one primary seed, one stability replication, and pair-clustered bootstrap intervals. Repeated measurements cannot compensate for the small designed corpus, author labels, correlated requested-seed outputs, infrastructure-failed trials, or the absence of independent failure domains.

\section{Defense Readiness}

\enlargethispage{2\baselineskip}
The repository is technically evidence-complete for the v0.7 rebuild: the complete test suite, Ruff, mypy, strengthened production branch-coverage gates, deterministic compiler benchmark, eight fixture-A2A replays, and the full three-model matrix pass their integrity checks. The 45-trial consensus-v4 matrix is complete but not accepted clean evidence because its provenance records the preserved dirty worktree. Defense readiness depends on presenting mixed results without dilution: RQ4 passes only safety superiority and symbolic invariants, while RQ5 is partially supported with two reconfiguration repetitions indeterminate. Independent corpus validation remains a disclosed limitation. Administrative tutor, tribunal, clean-source rerun, and packaging metadata remain release blockers. Independent label review and a broader low-overlap raw-language corpus are required before claiming a neutral benchmark, a superior planning pipeline, or a complete map of supported fault conditions.

Version 0.8 addresses the specific v0.7 semantic failure modes through evidence grounding, deterministic ambiguity, quotation and negation boundaries, non-authoritative retrieval, and constraint-directed provider search. Its complete development matrix passes the qualification gates for both permitted models, but this is intentionally a stopping point rather than a success claim: the same cases informed implementation revision. Qwen3-8B is excluded only from v0.8 and its v0.7 result remains unchanged. A cross-version comparison will be descriptive even after confirmation because the corpus and semantic interface changed. Until the two human labeling passes are completed and the 768 held-out observations are validated, v0.8 cannot alter the thesis's confirmatory conclusions.
